\documentclass[10pt,a4paper]{amsart}
\usepackage[margin=19mm]{geometry}
\usepackage{amsmath,amssymb,mathtools}
\usepackage[scr=rsfs,cal=euler]{mathalfa}
\usepackage{xfrac}
\usepackage[unicode,hidelinks,bookmarks=false]{hyperref}
\newcommand{\ie}{i.e.\ }
\newcommand{\cf}{cf.\ }
\newcommand{\resp}{respectively\ }
\newcommand{\Subsection}{\subsection}
\providecommand{\nobreakdash}{\nobreak\hbox{-}}
\setlength{\emergencystretch}{2em}
\multlinegap0pt
\usepackage{amssymb}
\usepackage{mathtools}
\usepackage[shortlabels]{enumitem}
\newtheorem{proposition}{Proposition}
\title{Almost Everywhere Convergence of Arithmetic Means of Walsh--Fourier Partial Sums Along Subsequences}
\author{Ushangi Goginava}
\date{Source: arXiv:2605.05146v1 (CC BY 4.0)}
\begin{document}
\maketitle

\noindent\emph{Source and licence.} Almost Everywhere Convergence of Arithmetic Means of Walsh--Fourier Partial Sums Along Subsequences. Version arXiv:2605.05146v1 (CC BY 4.0), distributed by arXiv under the Creative Commons Attribution 4.0 International licence, http://creativecommons.org/licenses/by/4.0/, as recorded in the arXiv licence metadata for this exact version and shown on the abstract page https://arxiv.org/abs/2605.05146v1. Full licence notice: \texttt{licenses/arxiv-2605.05146-CC-BY-4.0.txt}.\par\smallskip

\noindent\textit{Editorial scope.} Counted unit: the article's proposition prop:kernel-decomp (source0.tex lines 221--233). The article's complete original proof of it is delivered with it (source0.tex lines 235--267, one \texttt{proof} environment of 33 lines). No mathematics is written, repaired or re-derived: the statement and the proof are the article's own lines, in the article's own notation, and no retained line is substituted or re-flowed.  No further source-local context is needed: every cross-reference of every retained line resolves inside the two delivered spans.  Nothing was omitted from any retained span, so the package carries no omission record.  No retained line contains a citation, so the wrapper carries no bibliography and no bibliography entry.  No retained line contains a figure environment, an image-inclusion command or drawing code, so no image is delivered and the package declares no asset dependency.  Editorial formatting only, in addition: an AMS \texttt{amsart} preamble that restates the article's own declarations and macros used by the retained lines, namely the environments \texttt{proposition} on separate counters, together with the standard packages those lines need.  The retained environments print in this document as Proposition 1 in order of appearance, whereas the article prints the same items with the numbering of its own full document.  The complete native source of the article as distributed in the arXiv e-print for this exact version is bundled unchanged as \texttt{sources/199928/source0.tex}.
\begin{proposition}
\label{prop:kernel-decomp} For every $n\ge1$, 
\begin{equation}
D_n = D_{2^{|n|+1}}-\omega_n D_{2^{m(n)}}+\sum_{i=0}^{|n|-1}d_{n,i}.
\label{eq:kernel-decomp}
\end{equation}
Consequently, 
\begin{equation}
S_n f =\mathbb{E}_{|n|+1}f-\omega_n\mathbb{E}_{m(n)}(f\omega_n)+\widetilde
S_n f, \qquad \widetilde S_n f:=\sum_{i=0}^{|n|-1} f*d_{n,i}.
\label{eq:partial-sum-decomp}
\end{equation}
\end{proposition}

\begin{proof}
A standard identity for the Walsh Dirichlet kernel is 
\begin{equation}
D_n=\omega_n\sum_{i=0}^{|n|}n_i\bigl(D_{2^{i+1}}-D_{2^i}\bigr).
\label{eq:dir-standard}
\end{equation}
Since $n_i=0$ for $i<m(n)$, $n_{|n|}=1$, and $n_i=1-\lambda_{n,i}$ for $%
m(n)\le i<|n|$, we obtain 
\begin{equation*}
\sum_{i=0}^{|n|}n_i\bigl(D_{2^{i+1}}-D_{2^i}\bigr)
= \sum_{i=m(n)}^{|n|}\bigl(D_{2^{i+1}}-D_{2^i}\bigr)
- \sum_{i=0}^{|n|-1}\lambda_{n,i}\bigl(D_{2^{i+1}}-D_{2^i}\bigr). 
\end{equation*}
The first sum telescopes to $D_{2^{|n|+1}}-D_{2^{m(n)}}$. Multiplying by $%
\omega_n$ and using that $\omega_n\equiv1$ on the support of $D_{2^{|n|+1}}$%
, namely on $I_{|n|+1}$, gives 
\begin{equation*}
D_n = D_{2^{|n|+1}}-\omega_n D_{2^{m(n)}}
+\sum_{i=0}^{|n|-1}\lambda_{n,i}\,\omega_n\bigl(D_{2^i}-D_{2^{i+1}}\bigr), 
\end{equation*}
which is exactly \eqref{eq:kernel-decomp}.

Convolving \eqref{eq:kernel-decomp} with $f$ yields 
\begin{equation*}
S_n f= S_{2^{|n|+1}}f-(f*(\omega_nD_{2^{m(n)}}))+
\sum_{i=0}^{|n|-1}f*d_{n,i}. 
\end{equation*}
Now $S_{2^{|n|+1}}f=\mathbb{E}_{|n|+1}f$, and a standard Walsh character calculation shows that 
\begin{equation*}
(f*(\omega_nD_{2^{m(n)}}))(y)=\omega_n(y)\,\mathbb{E}_{m(n)}(f\omega_n)(y). 
\end{equation*}
Hence \eqref{eq:partial-sum-decomp} follows.
\end{proof}

\end{document}
